\begin{proofbox}

	\textbf{\textcolor{CProof}{思路提示}}

	利用不等式的性质：在 $ax>b$ 两边同时除以 $a$，但必须考虑 $a$ 的符号对不等号方向的影响；$a=0$ 时不能除，需单独判断。

	\medskip
	\textbf{\textcolor{CProof}{正式推导}}
	\begin{description}[style=nextline, leftmargin=2.8em, labelsep=0.5em, itemsep=0.55em]
		\item[\textbf{步骤一（正系数情形）：}] 若 $a>0$，则两边同除以 $a$，不等号不变：
			\[
				x > \frac{b}{a}.
			\]
			\par\textbf{理由：} 不等式性质：$a>0$ 时，$\frac{ax}{a} > \frac{b}{a}$ 等价于 $x > \frac{b}{a}$。
		\item[\textbf{步骤二（负系数情形）：}] 若 $a<0$，则两边同除以 $a$，不等号反向：
			\[
				x < \frac{b}{a}.
			\]
			\par\textbf{理由：} 不等式性质：$a<0$ 时，除以负数不等号转向，故得 $x<\frac{b}{a}$。
		\item[\textbf{步骤三（零系数退化）：}] 若 $a=0$，原不等式化为 $0 \cdot x > b$，即 $0 > b$。
			\par\textbf{理由：} 当 $a=0$ 时左边恒为 $0$，不等式等价于常数不等式 $0>b$。
		\item[\textbf{步骤四（零系数细分）：}] 对于 $a=0$ 的情况：
			\[
				\begin{cases}
					\mathbb{R}, & b<0,\\
					\varnothing, & b\ge 0.
			\end{cases}
			\]
			\par\textbf{理由：} 若 $b<0$，则 $0>b$ 恒成立，解集为全体实数；若 $b\ge 0$，$0>b$ 不可能成立，故无解。
	\end{description}

	\medskip
	\textbf{\textcolor{CProof}{结论回扣}}

	综上，对任意 $a,b\in\mathbb{R}$，不等式 $ax>b$ 的解集完全由 $a>0$、$a<0$、$a=0$ 三类情形决定，每一类都有确定的解集形式，从而得到完整的分类讨论规则。

\end{proofbox}
